% ==============================================================================
% CONCLUSION GÉNÉRALE
% ==============================================================================
\chapter*{Conclusion Générale}
\addcontentsline{toc}{chapter}{Conclusion Générale}
\markboth{CONCLUSION GÉNÉRALE}{}

\vspace{0.3cm}

Ce mémoire a présenté la conception et le développement d'un système
complet d'analyse automatique de sentiments appliqué aux commentaires
d'étudiants sur leurs cours universitaires. Ce projet représente une
intégration réussie de plusieurs domaines de l'informatique moderne :
l'intelligence artificielle, le développement web et mobile, les bases
de données relationnelles et les architectures distribuées.

\vspace{0.5cm}

\section*{Synthèse des réalisations}

\begin{tcolorbox}[
  enhanced,
  colback=bleuprincipal!5!white,
  colframe=bleuprincipal,
  boxrule=1.5pt, arc=6pt,
  left=10pt, right=10pt, top=8pt, bottom=8pt
]
\begin{enumerate}
  \item \textbf{Modèle d'IA :} Le modèle CamemBERT fine-tuné sur 225
        commentaires étudiants en français a atteint une précision de
        \textbf{79.41\%} sur le jeu de test, démontrant la puissance du
        transfer learning pour des tâches spécialisées avec peu de données.

  \item \textbf{API REST :} Le serveur FastAPI expose 5 endpoints documentés,
        gère la validation des données via Pydantic et persiste chaque analyse
        dans PostgreSQL via SQLAlchemy ORM. Le pattern Singleton garantit
        des temps de réponse inférieurs à 1 seconde.

  \item \textbf{Dashboard analytique :} Le tableau de bord Streamlit offre
        des visualisations interactives (camembert, histogrammes, jauges)
        permettant aux enseignants de suivre en temps réel les tendances
        de satisfaction par cours et par matière.

  \item \textbf{Application mobile :} L'application Flutter offre une
        interface intuitive, matérielle et moderne permettant aux étudiants
        de soumettre leurs commentaires, consulter les résultats en temps réel
        et naviguer dans l'historique depuis n'importe quel smartphone Android.
\end{enumerate}
\end{tcolorbox}

\vspace{0.5cm}

\section*{Bilan technique}

Ce projet a permis d'explorer et de maîtriser un large spectre de
technologies : PyTorch et HuggingFace Transformers pour le deep learning,
FastAPI et SQLAlchemy pour le backend, Flutter/Dart pour le mobile, et
Streamlit/Plotly pour la visualisation. Cette polyvalence technique
constitue un acquis précieux, directement transférable dans un contexte
professionnel.

Sur le plan algorithmique, la compréhension approfondie des mécanismes
d'attention des Transformers, du fine-tuning, de l'optimiseur AdamW et
du scheduler linéaire avec warmup a enrichi notre bagage en apprentissage
automatique moderne. La résolution des défis concrets — débogage du
pipeline de données, gestion des erreurs d'authentification PostgreSQL,
configuration de l'environnement Android — a également développé
des compétences pratiques essentielles.

\vspace{0.5cm}

\section*{Perspectives immédiates}

Les trois priorités d'amélioration identifiées sont :
\begin{itemize}
  \item \textbf{Augmentation des données :} collecter 500+ commentaires
        réels pour améliorer la précision au-delà de 85\%.
  \item \textbf{Sécurisation :} implémenter l'authentification JWT et le
        déploiement Docker pour une mise en production réelle.
  \item \textbf{Analyse d'aspects :} passer d'une classification globale
        à une analyse fine des aspects pédagogiques mentionnés.
\end{itemize}

\vspace{0.5cm}

\section*{Mot final}

Ce projet de fin d'études a été l'occasion de réaliser que la frontière
entre recherche et application concrète est plus mince qu'il n'y paraît.
CamemBERT, un modèle sorti de laboratoires de recherche de renom, peut
être adapté en quelques jours pour résoudre un problème réel et local.
Cette capacité à connecter la recherche fondamentale à des besoins
pratiques constitue, à notre sens, l'essence même de l'ingénierie
informatique moderne.

\vspace{0.5cm}

\begin{flushright}
\textit{``La technologie est à son meilleur quand elle rapproche les gens.''}\\
\vspace{0.3cm}
\textbf{Hissein Nasser} — Mai 2026
\end{flushright}

\cleardoublepage
